\documentclass[11pt]{article}
\usepackage[margin=0.85in]{geometry}
\usepackage{amsmath}
\usepackage{booktabs}
\usepackage{array}
\usepackage{float}

\title{Leveraged S\&P 500 Performance Analysis}
\author{Prepared by Harman Singh for Professor Phelim Boyle}
\date{June 2026}

\begin{document}
\maketitle

\section*{Data and Sample}
The analysis uses monthly S\&P 500 total return data from January 2017 to May 2026, for a total of 113 monthly observations.
The data source is the Yahoo Finance S\&P 500 Total Return Index series, \texttt{\string^SP500TR}, which includes dividends and is therefore consistent with the paper's dividend-reinvestment convention.
All statistics are reported monthly and are not annualised.

\section*{Methodology}
Let $r_t$ denote the S\&P 500 total return in month $t$, in percent, and let $r_f$ denote the constant monthly risk-free rate, also in percent.
For each leverage level $\lambda$, the portfolio return is
\[
R_t(\lambda)=\lambda r_t-(\lambda-1)r_f=\lambda(r_t-r_f)+r_f.
\]
The analysis uses $\lambda \in \{0.5,1,2,3,4,5,6\}$ and $r_f \in \{0.00\%,0.25\%,0.50\%\}$.

For each portfolio return series, the mean is
\[
\bar{R}=\frac{1}{n}\sum_{t=1}^n R_t,
\]
and the sample standard deviation is
\[
\sigma=\sqrt{\frac{1}{n-1}\sum_{t=1}^n (R_t-\bar{R})^2}.
\]
The Sharpe ratio is $\bar{R}/\sigma$.
The Sortino ratio uses the zero-return threshold:
\[
\sigma_d=\sqrt{\frac{1}{n}\sum_{t=1}^n \min(R_t,0)^2},
\qquad
\mathrm{Sortino}=\frac{\bar{R}}{\sigma_d}.
\]
The Omega ratio is
\[
\Omega=\frac{\sum_{t:R_t>0} R_t}{\sum_{t:R_t<0} |R_t|}.
\]

\section*{Results}
\begin{table}[H]
\centering
\caption{Leveraged S\&P 500 monthly statistics for $r_f=0.00\%$.}
\label{tab:rf-00}
\begin{tabular}{lrrrrr}
\toprule
$\lambda$ & $\bar{R}$ (\%) & $\sigma$ (\%) & Sharpe & Sortino & Omega \\
\midrule
0.5 & 0.6638 & 2.2719 & 0.2922 & 0.4734 & 2.0770 \\
1 & 1.3275 & 4.5437 & 0.2922 & 0.4734 & 2.0770 \\
2 & 2.6551 & 9.0875 & 0.2922 & 0.4734 & 2.0770 \\
3 & 3.9826 & 13.6312 & 0.2922 & 0.4734 & 2.0770 \\
4 & 5.3102 & 18.1750 & 0.2922 & 0.4734 & 2.0770 \\
5 & 6.6377 & 22.7187 & 0.2922 & 0.4734 & 2.0770 \\
6 & 7.9653 & 27.2624 & 0.2922 & 0.4734 & 2.0770 \\
\bottomrule
\end{tabular}
\end{table}

\begin{table}[H]
\centering
\caption{Leveraged S\&P 500 monthly statistics for $r_f=0.25\%$.}
\label{tab:rf-25}
\begin{tabular}{lrrrrr}
\toprule
$\lambda$ & $\bar{R}$ (\%) & $\sigma$ (\%) & Sharpe & Sortino & Omega \\
\midrule
0.5 & 0.7888 & 2.2719 & 0.3472 & 0.5852 & 2.3582 \\
1 & 1.3275 & 4.5437 & 0.2922 & 0.4734 & 2.0770 \\
2 & 2.4051 & 9.0875 & 0.2647 & 0.4205 & 1.9470 \\
3 & 3.4826 & 13.6312 & 0.2555 & 0.4033 & 1.9050 \\
4 & 4.5602 & 18.1750 & 0.2509 & 0.3948 & 1.8843 \\
5 & 5.6377 & 22.7187 & 0.2482 & 0.3897 & 1.8719 \\
6 & 6.7153 & 27.2624 & 0.2463 & 0.3863 & 1.8637 \\
\bottomrule
\end{tabular}
\end{table}

\begin{table}[H]
\centering
\caption{Leveraged S\&P 500 monthly statistics for $r_f=0.50\%$.}
\label{tab:rf-50}
\begin{tabular}{lrrrrr}
\toprule
$\lambda$ & $\bar{R}$ (\%) & $\sigma$ (\%) & Sharpe & Sortino & Omega \\
\midrule
0.5 & 0.9138 & 2.2719 & 0.4022 & 0.7058 & 2.6755 \\
1 & 1.3275 & 4.5437 & 0.2922 & 0.4734 & 2.0770 \\
2 & 2.1551 & 9.0875 & 0.2372 & 0.3696 & 1.8231 \\
3 & 2.9826 & 13.6312 & 0.2188 & 0.3367 & 1.7438 \\
4 & 3.8102 & 18.1750 & 0.2096 & 0.3205 & 1.7052 \\
5 & 4.6377 & 22.7187 & 0.2041 & 0.3109 & 1.6824 \\
6 & 5.4653 & 27.2624 & 0.2005 & 0.3045 & 1.6673 \\
\bottomrule
\end{tabular}
\end{table}

\section*{Brief Discussion}
When $r_f=0.00\%$, leverage scales both average return and volatility almost proportionally.
Sharpe, Sortino, and Omega are unchanged across leverage levels because every portfolio return is a constant multiple of the same underlying S\&P 500 total return series.

When $r_f=0.25\%$, borrowing costs reduce the return earned at higher leverage levels.
The mean return still rises with $\lambda$, but Sharpe, Sortino, and Omega decline as leverage increases because the cost of financing lowers returns without reducing downside exposure.

When $r_f=0.50\%$, the effect of borrowing costs is stronger.
The higher-leverage portfolios still have larger mean returns, but the risk-adjusted ratios fall more sharply, showing that leverage is less attractive when the monthly financing cost is higher.

\end{document}
